\documentclass[12pt,a4paper]{article}

% ----------------------------------------------------------------------
% Pacotes
% ----------------------------------------------------------------------
\usepackage[utf8]{inputenc}
\usepackage[T1]{fontenc}
\usepackage[brazil]{babel}
\usepackage[a4paper,margin=2.5cm]{geometry}
\usepackage{booktabs}
\usepackage{enumitem}
\usepackage{xcolor}
\usepackage{listings}
\usepackage{titlesec}
\usepackage[hidelinks]{hyperref}

% ----------------------------------------------------------------------
% Aparência
% ----------------------------------------------------------------------
\definecolor{azul}{RGB}{13,110,253}
\titleformat{\section}{\large\bfseries\color{azul}}{\thesection}{1em}{}
\titleformat{\subsection}{\normalsize\bfseries}{\thesubsection}{1em}{}
\setlength{\parindent}{0pt}
\setlength{\parskip}{6pt}

\lstset{
  basicstyle=\ttfamily\small,
  backgroundcolor=\color{gray!10},
  frame=single,
  rulecolor=\color{gray!40},
  breaklines=true,
  columns=fullflexible,
  showstringspaces=false
}

% ----------------------------------------------------------------------
% Capa / Título
% ----------------------------------------------------------------------
\title{\textbf{Sprint 5 --- Descrição do MVP Implementado}\\[4pt]
\large Salvou CRM}
\author{Engenharia de Software II}
\date{\today}

\begin{document}
\maketitle

\begin{flushleft}
\textbf{Projeto:} Salvou CRM \\
\textbf{Repositório:} \url{https://github.com/llagogabriel/salvou-crm} \\
\textbf{Branch:} \texttt{Verificações}
\end{flushleft}

\hrulefill

% ----------------------------------------------------------------------
\section{Visão geral do MVP}

O \textbf{Salvou CRM} é um sistema de gestão de relacionamento com clientes voltado a
profissionais e pequenos negócios de serviços (ex.: manicure, cabeleireiro, estética).
O objetivo do MVP é centralizar, em uma única aplicação web, o \textbf{cadastro de
clientes}, o \textbf{registro de vendas/serviços} e o \textbf{controle de uma agenda de
compromissos}, substituindo controles manuais (caderno, planilha) por uma ferramenta
organizada e com dados persistidos.

A aplicação foi construída em \textbf{Java~17} com o framework \textbf{Spring Boot~3.2},
seguindo o padrão \textbf{MVC (Model--View--Controller)}. A interface é renderizada no
servidor com \textbf{Thymeleaf} e \textbf{Bootstrap~5}, e os dados são persistidos em um
banco \textbf{H2 em arquivo} por meio do \textbf{Spring Data JPA / Hibernate}, garantindo
que as informações permaneçam salvas entre as execuções.

% ----------------------------------------------------------------------
\section{Funcionalidades implementadas}

\subsection{Módulo de Clientes}
\begin{itemize}[leftmargin=*]
  \item Cadastro de cliente com nome, telefone, CPF, endereço, e-mail e data de nascimento.
  \item \textbf{Regra de negócio:} não permite CPF duplicado (verificação via
        \texttt{existsByCpf}), exibindo mensagem de erro ao usuário.
  \item Listagem de todos os clientes cadastrados.
  \item Perfil do cliente, incluindo o cálculo do \textbf{valor total já gasto}
        (somatório de todas as vendas --- \texttt{getValorTotalVendas()}).
  \item Cadastro rápido de cliente via requisição assíncrona (endpoint REST
        \texttt{POST /clientes/rapido}), usado dentro da tela de vendas.
\end{itemize}

\subsection{Módulo de Vendas e Serviços}
\begin{itemize}[leftmargin=*]
  \item Registro de venda vinculando um \textbf{cliente} a um ou mais \textbf{serviços}.
  \item \textbf{Cálculo automático do valor total} a partir dos preços dos serviços
        (o servidor reconsulta os preços reais no banco, evitando manipulação pelo cliente).
  \item Data/hora da venda registrada automaticamente.
  \item Cadastro rápido de serviço (\texttt{POST /servicos/rapido}).
  \item Relacionamentos persistidos: \textbf{Cliente 1--N Venda} e \textbf{Venda N--N Serviço}.
\end{itemize}

\subsection{Módulo de Agenda}
\begin{itemize}[leftmargin=*]
  \item Cadastro de compromissos (atendimento, lembrete, pagamento, etc.), com descrição,
        tipo, data, observações e vínculo \textbf{opcional} a um cliente.
  \item Conclusão e exclusão de compromissos.
  \item Consultas filtradas: itens não concluídos e itens dos \textbf{próximos 7 dias}
        (\texttt{findByConcluidoFalse}, \texttt{findByDataAgendaBetween}).
\end{itemize}

\subsection{Tela inicial (Home)}
\begin{itemize}[leftmargin=*]
  \item Página de entrada (\texttt{/}) com menu de navegação para os módulos do sistema.
\end{itemize}

% ----------------------------------------------------------------------
\section{Fluxo completo: entrada $\rightarrow$ processamento $\rightarrow$ saída}

O sistema entrega o fluxo de ponta a ponta exigido pelo MVP. Exemplo do
\textbf{registro de uma venda}:

\begin{enumerate}[leftmargin=*]
  \item \textbf{Entrada:} o usuário acessa \texttt{/venda}, seleciona o cliente e os
        serviços no formulário (interface Thymeleaf/Bootstrap).
  \item \textbf{Processamento:} os dados são enviados ao \texttt{VendaController}, que busca
        o cliente e os serviços reais no banco, \textbf{calcula o valor total}, cria o
        vínculo entre as entidades e registra a data da venda.
  \item \textbf{Saída/Persistência:} a venda é salva no banco H2 via \texttt{VendaRepository},
        e o sistema retorna a confirmação. O valor passa a compor o \textbf{total gasto}
        exibido no perfil do cliente.
\end{enumerate}

O mesmo padrão (formulário $\rightarrow$ controller $\rightarrow$ repositório $\rightarrow$
banco $\rightarrow$ tela) se repete nos módulos de Clientes e Agenda, garantindo um ciclo
completo e coerente.

% ----------------------------------------------------------------------
\section{Tecnologias utilizadas}

\begin{center}
\begin{tabular}{ll}
\toprule
\textbf{Camada} & \textbf{Tecnologia} \\
\midrule
Linguagem              & Java 17 \\
Framework              & Spring Boot 3.2 \\
Web / Controllers      & Spring MVC \\
Visualização (View)    & Thymeleaf + Bootstrap 5 \\
Persistência           & Spring Data JPA + Hibernate \\
Banco de dados         & H2 (modo arquivo) \\
Build / Dependências   & Maven \\
\bottomrule
\end{tabular}
\end{center}

% ----------------------------------------------------------------------
\section{Como executar (para a demonstração)}

\begin{lstlisting}
# Na raiz do projeto (branch Verificacoes)
mvn spring-boot:run

# Acessar no navegador:
http://localhost:8080/
\end{lstlisting}

O banco H2 fica salvo em \texttt{./data/salvouDB}, preservando os dados cadastrados.
O console do banco está disponível em \texttt{http://localhost:8080/h2-console}
(caso seja necessário para evidência).

\end{document}
